
%(BEGIN_QUESTION)
% Copyright 2006, Tony R. Kuphaldt, released under the Creative Commons Attribution License (v 1.0)
% This means you may do almost anything with this work of mine, so long as you give me proper credit

Forklar virkemåten til disse aktuator-typene for ventiler, og svar deretter på spørsmålene:

\begin{itemize}
\item{} Fjær-og-membran (pneumatisk)
\item{} Stempel (pneumatisk)
\item{} Elektromekanisk (elektrisk motor)
%\item{} Elektromagnetisk (elektrisk solenoide)
%\item{} Elektrohydraulisk
\end{itemize}

\vskip 10pt

Hvis vi ser bort fra solenoideaktuator, som ikke brukes til proporsjonal regulering: hvilken aktuator-type i listen er enklest?

\vskip 10pt

Hvilke av aktuatorene i listen holder i seg selv posisjon ved bortfall av strøm eller trykkluftforsyning?

\vskip 10pt

Hvilke av aktuatorene i listen «feiler» i seg selv i én retning (enten helt åpen eller helt lukket) ved bortfall av forsyning, eller kan modifiseres til å gjøre det?

\underbar{file i00781}
%(END_QUESTION)





%(BEGIN_ANSWER)

Fjær-og-membran-aktuatorer virker ved at trykkluft presser mot en stor, fleksibel membran. Kraften fra lufttrykket på membranarealet virker mot kraften fra en komprimert fjær. Fjærkompresjonen må endres for å balansere nytt lufttrykk, og dette gir posisjonsendring som brukes til å bevege ventilmekanismen.

\vskip 10pt

Pneumatiske stempelaktuatorer bruker et stempel i stedet for membran. Noen stempelaktuatorer har fjær for returkraft (luft skyver stempelet én vei, fjæren skyver tilbake), mens andre er dobbeltvirkende (trykkluft på begge sider av stempelet).

\vskip 10pt

Elektromekaniske aktuatorer bruker elektrisk motor koblet til ventilspindel/aksel via gir. Når motoren roterer, endres ventilposisjonen. De brukes ofte til av/på-styring (helt åpen/helt lukket), men kan også utstyres med posisjonssensorer og servodrift for proporsjonal regulering.

\vskip 10pt

Solenoide (elektromagnetiske) ventilaktuatorer er enkle åpen/lukket-enheter, ikke for struperegulering. De bruker en elektromagnetspole til å trekke en ankerdel, ofte mot en fjærkraft.

\vskip 10pt

Elektrohydrauliske aktuatorer bruker normalt en elektrisk motor til å drive en hydraulikkpumpe, der væskestrømmen driver en stempelaktuator.

\vskip 10pt

Elektromekaniske aktuatorer (der ventil flyttes av en elektrisk motor) holder i seg selv posisjon ved forsyningsbortfall. Det samme gjelder elektrohydrauliske aktuatorer (der ventil flyttes av stempel drevet av hydraulikktrykk fra en pumpemotor).

\vskip 10pt

Alle aktuatorer med returfjær vil i seg selv «feile» til en konsistent posisjon ved forsyningsbortfall. Dette inkluderer pneumatiske membran- og stempelaktuatorer, samt solenoideaktuatorer.

%(END_ANSWER)





%(BEGIN_NOTES)

Pneumatiske fjær-og-membran-aktuatorer er ofte langt enklere enn aktuatorer som bruker elektriske motorer. Pneumatiske stempelaktuatorer krever ofte en spesialenhet kalt {\it positioner} for å konvertere lufttrykkssignal til entydig stempelposisjon, på grunn av friksjonseffekter i sylinderen. Med membran er aktuatorbevegelse derimot mer direkte funksjon av påført lufttrykk.

Dette betyr ikke at pneumatiske stempelaktuatorer {\it ikke kan} fungere uten positioner. Men siden mange ikke kan det, er membranaktuatorer enklere i design.

\vskip 10pt

Pneumatiske aktuatorer med fjærretur vil ofte gå til en ytterstilling («fail-closed» eller «fail-open») hvis luftforsyningen faller bort. Dobbeltvirkende stempelaktuatorer stopper ved tap av lufttrykk, men holder ikke nødvendigvis posisjon stivt mot krefter fra ventilsatsen. Fjærreturaktuatorer kan modifiseres til «hold-posisjon»-oppførsel ved å legge til en «locking valve» som isolerer membrankammeret når forsyningsluft faller bort.

\vskip 10pt

Elektrohydrauliske aktuatorer kan modifiseres til å feile i en forutsigbar retning ved strømbrudd ved å bruke returfjær i hydraulikksylinderen og en fail-open-solenoideventil som utjevner trykket i sylinderen ved strømtap.

%INDEX% Sluttkontrollelementer, ventil: aktuatortyper

%(END_NOTES)


